% !TeX program = pdflatex
\documentclass[10pt,a4paper]{article}

\usepackage[utf8]{inputenc}
\usepackage[T1]{fontenc}
\usepackage{lmodern}
\usepackage{geometry}
\usepackage[hidelinks]{hyperref}
\usepackage{xcolor}

\geometry{
  top=1.7cm,
  bottom=1.8cm,
  left=1.8cm,
  right=1.8cm
}

\pagestyle{empty}
\frenchspacing
\setlength{\parindent}{0pt}
\setlength{\parskip}{0.45em}

\definecolor{mydarkgreen}{RGB}{0,0,0}

\begin{document}

\begin{minipage}[t]{0.48\textwidth}
\textbf{\large Javier Andres TARAZONA JIMENEZ}\\
Étudiant ingénieur IA \\
Machine Learning, Vision 3D, LLMs et prototypage \\
Massy, Île-de-France, France\\
\href{mailto:javitar06@gmail.com}{javitar06@gmail.com}\\
+33 07 46 36 97 75\\
\href{https://www.linkedin.com/in/javier-andres-tarazona-jimenez-84b489222/}{LinkedIn}
\end{minipage}
\hfill
\begin{minipage}[t]{0.42\textwidth}
\raggedleft
MBDA France \\
À l’attention de l’équipe Recrutement / Direction Technical \\
Site du Plessis-Robinson \\
1 Av. Réaumur \\
92350 Le Plessis-Robinson \\
France \\
\end{minipage}

\vspace{1.1cm}

\raggedleft
Palaiseau, le 24 juin 2026

\vspace{0.7cm}

\raggedright
\textbf{Objet : Candidature à l’apprentissage Ingénieur Data Scientist en Détection d’Anomalies F/H – Réf. R34045 – disponibilité 1 an}

\vspace{0.5cm}

{\small

Madame, Monsieur,

Actuellement étudiant colombien en double diplôme entre l’Universidad Nacional de Colombia et ENSTA Paris, au sein de l’Institut Polytechnique de Paris, je souhaite vous présenter ma candidature pour le poste d’Apprentissage Ingénieur Data Scientist en Détection d’Anomalies F/H, référence R34045, à partir de septembre 2026. J’ai bien noté que l’offre est prévue pour une durée de deux ans ; dans mon cas, je recherche une alternance d’un an dans le cadre de ma dernière année de cycle ingénieur. Je me permets néanmoins de vous transmettre ma candidature, au cas où mon profil pourrait être utile à MBDA, que ce soit pour cette mission ou pour une autre opportunité proche de vos besoins.

J’ai choisi l’informatique par besoin de comprendre les systèmes complexes, par goût pour la construction de solutions concrètes à partir d’idées abstraites, par recherche d’autonomie technique et par ambition académique. L’intelligence artificielle s’est ensuite imposée comme une évolution naturelle de ce parcours : passer de la programmation de systèmes à la conception de systèmes capables de percevoir, d’apprendre et d’aider à la décision. Je suis particulièrement attiré par les problèmes ouverts, ambigus et exigeants, où il faut articuler raisonnement, modélisation, expérimentation et qualité logicielle. La vision par ordinateur m’a permis de relier ce goût pour l’abstraction à une dimension très appliquée : transformer des mathématiques et du code en perception du monde réel. Plus largement, l’IA m’intéresse parce qu’elle constitue une technologie transversale, capable d’avoir un impact visible sur l’industrie, la société et l’économie.

Je suis intéressé par MBDA car votre environnement technique correspond à l’usage de l’intelligence artificielle qui m’attire le plus : une IA appliquée à des systèmes complexes, à la simulation numérique et à des problématiques industrielles exigeantes. La mission de détection d’anomalies sur des résultats de simulation me semble particulièrement cohérente avec mon parcours en IA, vision 3D, génération de données synthétiques et développement Python/PyTorch. Elle me permettrait de mettre mes compétences au service d’outils robustes, validés et intégrés dans un contexte opérationnel à forte valeur technologique.

Cette alternance représenterait pour moi une étape importante, car elle me permettrait de sortir d’une approche expérimentale de l’IA pour travailler sur une IA industrielle : structurée autour des données, de la validation, des métriques, de l’état de l’art, de l’amélioration continue et de l’intégration dans un environnement technique réel. Ce qui m’intéresse particulièrement n’est pas seulement d’entraîner un modèle performant, mais de comprendre comment rendre un modèle utile, justifiable, maintenable et intégrable pour des ingénieurs travaillant sur des systèmes critiques. Dans cette perspective, MBDA pourrait m’aider à transformer mon profil d’étudiant fortement orienté IA et vision par ordinateur en profil d’ingénieur capable de développer, valider et industrialiser des modèles pour des systèmes complexes.

Je pense pouvoir apporter à MBDA une combinaison cohérente de compétences en intelligence artificielle, simulation, vision par ordinateur et développement logiciel. À ENSTA Paris – U2IS, je travaille sur la génération de données 3D synthétiques avec Python/Blender, le rendu multi-vues, la structuration de métadonnées et l’entraînement de modèles PyTorch et Vision Transformers pour apprendre la structure 3D par prédiction de vues. Chez Engin A.I, j’ai développé et refactorisé des pipelines Python/OpenCV pour l’analyse d’images routières, avec un effort important sur la modularité, la performance et la maintenabilité. Mes projets, notamment Tameo et Slow, m’ont également permis de travailler sur des pipelines de vision, de détection, de suivi et d’évaluation de modèles. Enfin, ma formation à ENSTA et à l’UNAL m’a donné des bases solides en machine learning, estimation statistique, modélisation, simulation, optimisation, génie logiciel, bases de données et systèmes distribués.

Motivé par les défis technologiques de l’industrie de défense et par l’application rigoureuse de l’IA à des systèmes à forte responsabilité, je serais heureux d’échanger avec vous sur l’adéquation de mon profil avec vos besoins.

Je vous remercie pour l’attention portée à ma candidature et vous prie d’agréer, Madame, Monsieur, l’expression de mes salutations distinguées.

\textbf{Cordialement},

Javier Andres TARAZONA JIMENEZ


}

\end{document}